% !TEX program = xelatex
\documentclass[11pt,a4paper]{article}

% ---------------------------------------------------------------------------
% Encoding & Fonts  (XeLaTeX)
% ---------------------------------------------------------------------------
\usepackage{fontspec}
\usepackage{xeCJK}
\setCJKmainfont{PingFang TC}
\setmainfont{Times New Roman}

% ---------------------------------------------------------------------------
% Layout
% ---------------------------------------------------------------------------
\usepackage[top=2.2cm,bottom=2.2cm,left=2.5cm,right=2.5cm]{geometry}
\usepackage{setspace}
\usepackage{parskip}
\setstretch{1.2}
\setlength{\parskip}{4pt}

% ---------------------------------------------------------------------------
% Section titles
% ---------------------------------------------------------------------------
\usepackage{titlesec}
\usepackage{xcolor}

\definecolor{ntublue}{RGB}{26,26,94}
\definecolor{lightblue}{RGB}{235,238,250}
\definecolor{lightgreen}{RGB}{241,248,241}
\definecolor{lightorange}{RGB}{255,245,236}
\definecolor{colegray}{RGB}{240,240,248}
\definecolor{warnorange}{RGB}{196,92,0}
\definecolor{successgreen}{RGB}{46,125,50}

\titleformat{\section}
  {\large\bfseries\color{ntublue}}{\thesection.}{0.6em}{}
  [\color{ntublue}\titlerule]
\titleformat{\subsection}
  {\normalsize\bfseries\color{ntublue}}{\thesubsection}{0.5em}{}
\titlespacing{\section}{0pt}{12pt}{4pt}
\titlespacing{\subsection}{0pt}{8pt}{3pt}

% ---------------------------------------------------------------------------
% Headers & Footers
% ---------------------------------------------------------------------------
\usepackage{fancyhdr}
\pagestyle{fancy}
\fancyhf{}
\renewcommand{\headrulewidth}{0.4pt}
\fancyhead[L]{\small\textcolor{ntublue}{Yahoo Shopping Seller Price Monitor}}
\fancyhead[R]{\small\textcolor{ntublue}{BDS Final Project --- Spring 2026, NTU}}
\fancyfoot[C]{\small\thepage}

% ---------------------------------------------------------------------------
% Tables
% ---------------------------------------------------------------------------
\usepackage{booktabs}
\usepackage{tabularx}
\usepackage{array}
\usepackage{ragged2e}
\newcolumntype{L}[1]{>{\RaggedRight\arraybackslash}p{#1}}
\newcolumntype{C}[1]{>{\Centering\arraybackslash}p{#1}}

% ---------------------------------------------------------------------------
% Callout boxes
% ---------------------------------------------------------------------------
\usepackage{tcolorbox}
\tcbuselibrary{skins,breakable}

\newtcolorbox{insightbox}{
  breakable,colback=lightblue,colframe=ntublue,
  leftrule=4pt,toprule=0pt,bottomrule=0pt,rightrule=0pt,
  arc=0pt,left=6pt,right=6pt,top=4pt,bottom=4pt,
}
\newtcolorbox{successbox}{
  breakable,colback=lightgreen,colframe=successgreen,
  leftrule=4pt,toprule=0pt,bottomrule=0pt,rightrule=0pt,
  arc=0pt,left=6pt,right=6pt,top=4pt,bottom=4pt,
}
\newtcolorbox{warnbox}{
  breakable,colback=lightorange,colframe=warnorange,
  leftrule=4pt,toprule=0pt,bottomrule=0pt,rightrule=0pt,
  arc=0pt,left=6pt,right=6pt,top=4pt,bottom=4pt,
}

% ---------------------------------------------------------------------------
% Code blocks
% ---------------------------------------------------------------------------
\usepackage{listings}
\lstset{
  basicstyle=\ttfamily\scriptsize,
  backgroundcolor=\color{colegray},
  frame=single,framerule=0.4pt,rulecolor=\color{gray!50},
  breaklines=true,keepspaces=true,
  showspaces=false,showstringspaces=false,
}

% ---------------------------------------------------------------------------
% Misc
% ---------------------------------------------------------------------------
\usepackage{enumitem}
\setlist[itemize]{leftmargin=1.4em,itemsep=1pt,parsep=0pt,topsep=2pt}
\setlist[enumerate]{leftmargin=1.4em,itemsep=1pt,parsep=0pt,topsep=2pt}
\usepackage{float}
\usepackage{caption}
\captionsetup{font=small,labelfont=bf,skip=4pt}
\usepackage{hyperref}
\hypersetup{
  colorlinks=true,linkcolor=ntublue,urlcolor=ntublue,
  pdftitle={Yahoo Shopping Seller Competitor Price Monitoring Service},
}
\usepackage{graphicx}

% ===========================================================================
\begin{document}
% ===========================================================================

% ---------------------------------------------------------------------------
% Cover Page
% ---------------------------------------------------------------------------
\thispagestyle{empty}
\begin{center}
  \vspace*{2.5cm}
  {\LARGE\bfseries\color{ntublue}
    Seller Competitor Price Monitoring Service}\\[0.5em]
  {\large Built on Yahoo Shopping Taiwan}\\[0.3em]
  {\normalsize Designing a System That Monetizes Data}\\[1.5em]
  \rule{0.55\textwidth}{0.8pt}\\[1.2em]
  {\normalsize Big Data Systems (BDS) --- Final Project \quad|\quad
   Spring 2026 \quad|\quad National Taiwan University}\\[2.5cm]

  \begin{tabular}{rl}
    \textbf{Student ID:} & \texttt{<your-student-id>} \\[0.3em]
    \textbf{GitHub:}     & \url{https://github.com/<your-id>/bds-price-monitor} \\[0.3em]
    \textbf{Live Demo:}  & \url{https://bds-price-monitor.railway.app} \\[0.3em]
    \textbf{Date:}       & June 2026 \\
  \end{tabular}
\end{center}
\newpage

% ===========================================================================
\section{Target Customer}
% ===========================================================================

\subsection{Segment and Persona}

Our customers are \textbf{small-to-mid-sized sellers on Taiwan e-commerce platforms}
--- sole proprietors or micro-businesses selling 50--500 SKUs (electronics
accessories, fashion, household goods) with monthly GMV of NT\$50K--500K and
net margins of 10--15\%. We focus on sellers who list products on
\textbf{Yahoo Shopping Taiwan} (\texttt{tw.buy.yahoo.com}), one of Taiwan's
largest e-commerce marketplaces.

\textbf{Representative persona --- Wang Wei, Yahoo Shopping seller since 2020:}
he sells approximately 200 electronics-accessory listings at NT\$150K GMV/month.
His core pain point is discovering that a competitor slashed prices only
\emph{after} his own sales drop and his platform ranking has already fallen.
His current workaround is manual morning checks of competitor listings,
which means he misses all off-hours flash sales and has no access to price history.

\subsection{Why Our System Wins}

\begin{table}[H]
\centering\small\renewcommand{\arraystretch}{1.25}
\begin{tabularx}{\textwidth}{X X}
\toprule
\textbf{Status Quo} & \textbf{Our System} \\
\midrule
Manual checks 1--2$\times$/day; misses off-hours
  & Automated 30-minute scrape cycle, 24/7 \\
No price-history data
  & 30-day trend chart per competitor product \\
No proactive alerts
  & Email alert within minutes of threshold breach \\
Seller must manually find competitor links
  & Automatic competitor discovery from product name \\
English-only global tools (Prisync, Price2Spy)
  & Built for Yahoo Shopping TW; NT\$ pricing \\
Prisync entry: \$9\,USD/mo ($\approx$NT\$288)
  & NT\$149/mo --- roughly half the cost of Prisync, locally optimised \\
\bottomrule
\end{tabularx}
\caption{Status quo vs.\ our system}
\end{table}

\textbf{Addressable market:} Yahoo Shopping Taiwan hosts tens of thousands of
active third-party sellers. Targeting the tool-buying top 10\% at NT\$149/month
yields a \textbf{NT\$7M+/year beachhead} market with clear expansion paths
to Momo and PChome.

% ===========================================================================
\section{Evidence of Demand and Willingness to Pay}
% ===========================================================================

We gathered evidence from \textbf{four independent public data sources}
using Python scripts in \texttt{demand\_analysis/} (run
\texttt{python run\_all.py} to reproduce all figures).

\begin{table}[H]
\centering\small\renewcommand{\arraystretch}{1.25}
\begin{tabularx}{\textwidth}{L{3.2cm} L{3.4cm} C{1.6cm} X}
\toprule
\textbf{Source} & \textbf{Method} & \textbf{Result} & \textbf{Signal} \\
\midrule
PTT BBS (e-shopping, BuyTogether)
  & PTT search API; regex parse
  & 58 threads (40 = 69\% with demand signal)
  & Sellers complain about rising fees, margin erosion, and competitor dumping \\
Google News RSS (zh-TW)
  & RSS feed; headline de-duplication
  & 163 unique articles
  & Mainstream press repeatedly covers ``seller profits squeezed'' \\
App Store reviews (iTunes RSS)
  & Keyword match on latest 50 reviews
  & 6 / 50 = 12\% hit rate
  & Shopee Seller Center app rated \textbf{1.6\,$\bigstar$} vs.\ 4.7 for buyer app --- severe unmet seller needs \\
iTunes Lookup API
  & REST JSON; rating count
  & 780,163 buyer ratings; 50 seller ratings
  & Large buyer base confirms market scale; tiny seller-app engagement confirms tooling gap \\
\bottomrule
\end{tabularx}
\caption{Demand evidence: four independent public data sources}
\end{table}

\begin{insightbox}
The 15,600$\times$ gap in App Store ratings between buyer apps (780K ratings)
and seller management tools (50 ratings) reflects chronic under-investment in
seller tooling across Taiwan e-commerce, creating a clear opening for a focused
third-party monitoring service.
\end{insightbox}

\subsection{Competitor Pricing Benchmark and Willingness to Pay}

\begin{table}[H]
\centering\small\renewcommand{\arraystretch}{1.25}
\begin{tabularx}{\textwidth}{L{2.6cm} L{3cm} C{2.4cm} C{2cm} C{2cm}}
\toprule
\textbf{Tool} & \textbf{Focus} & \textbf{Entry Price} & \textbf{Language} & \textbf{Yahoo TW?} \\
\midrule
Prisync   & General e-commerce & \$9\,USD/mo ($\approx$NT\$288)  & English & No \\
Price2Spy & General e-commerce & \$19\,USD/mo ($\approx$NT\$608) & English & No \\
Yahoo built-in analytics & Yahoo only & Free & Chinese & Partial \\
\textbf{Our System} & \textbf{Yahoo Shopping TW} & \textbf{NT\$149/mo} & \textbf{English + Chinese} & \textbf{Yes} \\
\bottomrule
\end{tabularx}
\caption{Competitor pricing benchmark}
\end{table}

\begin{successbox}
\textbf{WTP conclusion:} Existing tools start at NT\$288--608/month but require
English proficiency and lack Taiwan-specific platform integration. Our NT\$149/month
entry is roughly half the price of the cheapest English competitor (Prisync at
NT\$288) while delivering automatic competitor discovery out of the box.
Sellers' demonstrated willingness to pay for platform ads and third-party ERP
tools (NT\$500--2,000/month) confirms this segment will pay for margin-protecting
tools even at higher price points.
\end{successbox}

\subsection{Pricing Rationale}

We set the entry price at \textbf{NT\$149/month} based on three considerations:

\begin{enumerate}
  \item \textbf{Accessibility for small sellers.}
        Our smallest target seller generates NT\$50K GMV/month at a 10\% margin,
        yielding roughly NT\$5,000/month in gross profit.
        NT\$149 represents less than \textbf{3\% of that margin},
        making the subscription an easy expense to justify without managerial
        approval --- comparable to a single platform ad campaign.

  \item \textbf{Undercut the cheapest international alternative.}
        Prisync, the lowest-priced global competitor, costs \$9\,USD/month
        ($\approx$NT\$288). Pricing at NT\$149 --- just over half --- gives
        price-sensitive Taiwan sellers a compelling reason to choose a
        locally focused tool even before evaluating features.
        Price2Spy at NT\$608 makes our price look even more attractive.

  \item \textbf{Infrastructure sustainability.}
        At NT\$149/month, the system covers its full cloud infrastructure cost
        (Railway + Upstash, $\approx$NT\$800/month) at just 6 subscribers,
        and breaks even on one developer's salary at roughly 410 users ---
        approximately 10\% of the addressable beachhead.
        This is a realistic early-adoption target achievable through community
        word-of-mouth without paid advertising.
\end{enumerate}

% ===========================================================================
\newpage
\section{System Design}
% ===========================================================================

\subsection{Platform Selection: From Shopee to Yahoo Shopping}

Our initial prototype targeted \textbf{Shopee Taiwan}, which has a large seller
base and public product listings. During implementation, however, we discovered
that Shopee's v4 API applies a multi-layer anti-bot defence that makes
programmatic access infeasible:

\begin{warnbox}
\textbf{Shopee Anti-Bot Barrier.}
Even with a valid logged-in session cookie, all API calls to
\texttt{shopee.tw/api/v4/item/get} return error code \texttt{90309999}
(\texttt{action\_type: 2}). The server demands dynamically generated
\texttt{x-sap-access-*} authentication headers that are computed
by obfuscated JavaScript running inside the Shopee web application.
These tokens change per request and cannot be replicated without executing
a full browser environment (e.g.\ Selenium or Playwright),
which is 10--100$\times$ more resource-intensive and brittle at the
scale of a 30-minute polling cycle across many products.
\end{warnbox}

We therefore pivoted to \textbf{Yahoo Shopping Taiwan}
(\texttt{tw.buy.yahoo.com}), which imposes no equivalent barrier:

\begin{itemize}
  \item A simple HTTP \texttt{GET} to
        \texttt{https://tw.buy.yahoo.com/gdsale/gdsale.asp?gdid=\{GDID\}}
        returns a full HTML page with the product price embedded as a JSON
        literal (e.g.\ \texttt{"price": 440}).
  \item No authentication, no dynamic token, and no JavaScript execution
        are required. Standard browser \texttt{User-Agent} and
        \texttt{Referer} headers are sufficient.
  \item The search endpoint
        \texttt{https://tw.buy.yahoo.com/search/product?p=\{keyword\}}
        returns an HTML page containing \texttt{gdid} values for matching
        products, enabling automatic competitor discovery.
  \item Behaviour verified on 20+ real products across electronics, fashion,
        and home goods categories.
\end{itemize}

\subsection{Architecture}

\begin{figure}[H]
\begin{lstlisting}[language={}]
Yahoo Shopping Taiwan  (tw.buy.yahoo.com)
      |  Plain HTTP GET -- no auth, no token required
      |  (requests spider + APScheduler, every 30 min)
      v
[Kafka]  topic: price_events          (confluentinc/cp-kafka:7.6.1)
      |
      v
[Spark Streaming]  PySpark 3.5.1 -- foreachBatch:
   1. Read latest_price:{product_id} from Redis cache
   2. Compute price-change %
   3. If competitor price < seller my_price by >= alert threshold
         -> write Alert record to PostgreSQL
         -> send Gmail SMTP alert email to seller
   4. Update Redis  (TTL 3600 s)
   5. Append row to price_history in PostgreSQL
      |
      +---------------------+
      v                     v
[PostgreSQL 15]         [Redis 7]
 users                   latest_price:{product_id}
 products
   (parent + competitors,
    self-referential FK)
 price_history
 alerts
      |
      v
[FastAPI :8000]  <-->  [Flask Dashboard :5001]
 JWT auth REST API        Chart.js 30-day price-history charts
 Yahoo spider (on-demand) Competitor price comparison table
 Auto competitor discovery Email alert history log
\end{lstlisting}
\caption{End-to-end system architecture (8 Docker Compose services)}
\end{figure}

\subsection{Automatic Competitor Discovery}

A key differentiator of our system is \textbf{zero-configuration competitor
discovery}. When a seller adds their product URL, the system automatically
builds a competitor watchlist without any additional user input:

\begin{enumerate}
  \item \textbf{Keyword extraction.} The product name is tokenised and cleaned:
        English brand tokens (e.g.\ ``Esense'', ``Anker'') and any immediately
        following Chinese brand names are skipped; material/spec descriptors
        (e.g.\ aluminum alloy, silicone) are filtered out; remaining
        product-type acronyms (USB, HUB, HDMI, SSD, \ldots) and Chinese
        category words are combined into a 2--4 token search query.
        Example: ``Esense S649 A+C Aluminum USB3.2 HUB Hub'' $\to$ \texttt{USB HUB Hub}.
  \item \textbf{Yahoo search.} The keyword is submitted to the Yahoo Shopping
        search endpoint and up to 5 distinct \texttt{gdid} product IDs are
        extracted from the results page via regex.
  \item \textbf{Competitor fetch.} Each candidate \texttt{gdid} is fetched
        individually; product name and current price are extracted from the
        embedded HTML JSON.
  \item \textbf{Database insertion.} Each competitor is saved with
        \texttt{competitor\_of = parent\_id}, creating a self-referential
        relationship inside the \texttt{products} table.
        A PostgreSQL \texttt{SAVEPOINT} is used per competitor so that a
        single failure does not roll back the others.
  \item \textbf{Fallback re-discovery.} If a parent product has zero active
        competitors (e.g.\ due to a transient network timeout during the
        original add), triggering ``Scrape All Now'' automatically re-runs
        competitor discovery before the regular price-scrape cycle.
\end{enumerate}

\begin{successbox}
\textbf{Result in testing:} Adding a USB hub product at NT\$440 automatically
populated 5 competitors --- YUCUN (NT\$588), Esense alternate model (NT\$297),
ANTIAN (NT\$760 and NT\$1,080), and INTOPIC (NT\$425) --- within approximately
45 seconds. The NT\$297 competitor immediately triggered an email alert
(32.5\% cheaper than the seller's listed price).
\end{successbox}

\subsection{Technology Choices}

\begin{table}[H]
\centering\small\renewcommand{\arraystretch}{1.25}
\begin{tabularx}{\textwidth}{L{2.8cm} L{2.8cm} X}
\toprule
\textbf{Technology} & \textbf{Role} & \textbf{Justification} \\
\midrule
Apache Kafka (Confluent 7.6) & Message queue
  & Decouples scraper from processor; provides durable replay; supports
    future consumers such as ML anomaly detection or multi-user fan-out \\
Spark Streaming (PySpark 3.5) & Micro-batch processor
  & \texttt{foreachBatch} enables flexible Python logic for per-product
    threshold comparison; scales to thousands of concurrent products \\
PostgreSQL 15 & Persistent store
  & ACID guarantees; self-referential FK for parent/competitor products;
    composite index on \texttt{(product\_id, scraped\_at DESC)} for fast
    time-series queries \\
Redis 7 & Latest-price cache
  & O(1) lookup avoids a PostgreSQL round-trip on every Spark micro-batch;
    TTL = 3600\,s prevents stale cached values \\
Python \texttt{requests} + regex & Yahoo scraper
  & Lightweight HTTP scraping; Yahoo pages require no JavaScript execution;
    price extracted from embedded JSON with a single regex pattern \\
FastAPI + Flask & API + Dashboard
  & JWT-authenticated REST backend; on-demand scrape endpoints; Chart.js
    30-day price-history charts; Gmail SMTP alert emails \\
\bottomrule
\end{tabularx}
\caption{Technology choices and justifications}
\end{table}

\subsection{Data Schema}

The database contains four tables (see \texttt{init-db/schema.sql}):

\begin{itemize}
  \item \textbf{users}: \texttt{id, email, password\_hash, plan, created\_at}
  \item \textbf{products}: \texttt{id, user\_id, name, yahoo\_gdid, my\_price,
        alert\_threshold\_pct, active, competitor\_of, created\_at}\\
        The \texttt{competitor\_of} column is a \emph{self-referential} foreign
        key (\texttt{products.id $\to$ products.id ON DELETE CASCADE}).
        Parent products (added by the seller) have \texttt{competitor\_of = NULL};
        auto-discovered competitors have \texttt{competitor\_of = parent\_id}.
        A \texttt{UNIQUE(user\_id, yahoo\_gdid)} constraint prevents duplicate
        tracking of the same product by the same user.
  \item \textbf{price\_history}: \texttt{id, product\_id, price, stock,
        sold\_count, scraped\_at} with a composite index on
        \texttt{(product\_id, scraped\_at DESC)}
  \item \textbf{alerts}: \texttt{id, product\_id, old\_price, new\_price,
        change\_pct, email\_sent, created\_at}
\end{itemize}

\begin{lstlisting}[language=SQL, caption={Key products table definition}]
CREATE TABLE products (
    id                  SERIAL PRIMARY KEY,
    user_id             INTEGER REFERENCES users(id) ON DELETE CASCADE,
    name                VARCHAR(500) NOT NULL,
    yahoo_gdid          BIGINT NOT NULL,
    my_price            NUMERIC(10,2),
    alert_threshold_pct NUMERIC(5,2) DEFAULT 5.0,
    active              BOOLEAN DEFAULT TRUE,
    competitor_of       INTEGER REFERENCES products(id) ON DELETE CASCADE,
    created_at          TIMESTAMP DEFAULT NOW(),
    UNIQUE(user_id, yahoo_gdid)
);
\end{lstlisting}

\subsection{Scalability}

\begin{table}[H]
\centering\small\renewcommand{\arraystretch}{1.2}
\begin{tabularx}{\textwidth}{C{1.6cm} C{2.2cm} L{3.2cm} X}
\toprule
\textbf{Scale} & \textbf{Products} & \textbf{Bottleneck} & \textbf{Solution} \\
\midrule
MVP    & $\sim$100   & Single Kafka broker   & Docker Compose on one machine \\
10$\times$  & $\sim$1,000  & PostgreSQL write throughput & PgBouncer connection pool; monthly table partitioning \\
100$\times$ & $\sim$10,000 & Scraper concurrency vs.\ rate limits & Async \texttt{httpx}; per-domain request throttle \\
\bottomrule
\end{tabularx}
\caption{Scalability at different growth stages}
\end{table}

\subsection{Deployment}

Full system: \textbf{Docker Compose} (8 services ---
\texttt{zookeeper}, \texttt{kafka}, \texttt{postgresql}, \texttt{redis},
\texttt{scraper}, \texttt{pipeline}, \texttt{api}, \texttt{dashboard}),
launched with a single \texttt{docker-compose up --build}. For cloud
deployment, stateless services run on \textbf{Railway} (with managed
PostgreSQL and Redis plugins), and the local Kafka broker is replaced by
\textbf{Upstash} managed Kafka.

% ===========================================================================
\section{Bonus --- Go-to-Market Difficulties}
% ===========================================================================

\textbf{1.\ Why Not Shopee?}

\begin{warnbox}
We originally built the scraper targeting Shopee Taiwan and encountered an
insurmountable technical barrier: Shopee's v4 API requires dynamically generated
\texttt{x-sap-access-*} authentication tokens computed by obfuscated in-browser
JavaScript. These tokens change per request and cannot be obtained without running
a full headless browser. Even with a valid logged-in session cookie, all requests
return error \texttt{90309999}. Running Selenium or Playwright at scale is
10--100$\times$ more resource-intensive and brittle, making it impractical for a
30-minute polling service. Yahoo Shopping Taiwan imposes no equivalent barrier,
making it the viable platform for an MVP built with standard Python \texttt{requests}.
\end{warnbox}

\textbf{2.\ Legal Risk.}
Yahoo Shopping Taiwan's Terms of Service restrict systematic automated access
for commercial purposes. Our system scrapes publicly visible product pages at a
polite rate (30-minute intervals with 1.5-second inter-request delays),
consistent with how a human would browse. Mitigations include:
(a) pursuing a Yahoo partner API arrangement once the business scales;
(b) aggressive caching to minimise request volume;
(c) transparent disclosure of the scraping mechanism to users.

\textbf{3.\ Trust and Adoption.}
Our system only requires a competitor's product URL --- not store credentials ---
minimising perceived privacy risk. The primary go-to-market channel is Taiwan
seller Facebook communities and forums, where peer recommendations carry outsized
influence.

\textbf{4.\ Cold-Start.}
There is no network-effect dependency: a seller with a single product URL receives
immediate value. The automatic competitor-discovery feature (Section~3.3) means
sellers do not need to manually find competitor links; the system generates a
fully populated watchlist from the very first product added.

\textbf{5.\ Competitive Landscape.}

\begin{table}[H]
\centering\small\renewcommand{\arraystretch}{1.2}
\begin{tabularx}{\textwidth}{L{3.8cm} C{2.2cm} X}
\toprule
\textbf{Competitor} & \textbf{Threat} & \textbf{Our Moat} \\
\midrule
Yahoo built-in analytics & High & Shows only own sales data, not competitor prices or history \\
Prisync / Price2Spy      & Low  & No Yahoo TW URL parser; English-only; no auto-discovery \\
Developer copy           & Medium & First-mover brand in seller community; data network effect \\
\bottomrule
\end{tabularx}
\caption{Competitive landscape}
\end{table}

Long-term moat: \emph{data network effects} --- aggregated anonymised pricing
trends across thousands of sellers become a premium product
(e.g.\ ``average price movement in your category this week'').

\textbf{6.\ Unit Economics.}
Revenue: NT\$149/user/month.
Infrastructure (Railway + Upstash Kafka) for 100 users: approximately NT\$800/month.
Break-even on infrastructure: \textbf{6 users}.
Cost to sustain one developer (NT\$60K/month): \textbf{$\approx$410 users} ($\approx$10\% of beachhead).
At 20\% penetration ($\approx$NT\$1.2M/month) the system sustains a small team
and funds multi-platform expansion to Momo and PChome.

\end{document}
